\section{Registro de riesgos}
\label{sec:diseno-riesgos}

La siguiente tabla es el esqueleto del registro de riesgos. Las filas marcadas
como pendientes son ejemplos de marcador de posición y deben sustituirse por
riesgos reales del proyecto.

\begin{table}[htbp]
  \centering
  \caption{Registro de riesgos (esqueleto pendiente de completar).}
  \label{tab:registro-riesgos}
  \footnotesize
  \setlength{\tabcolsep}{4pt}
  \begin{tabularx}{\textwidth}{@{}l *{4}{>{\raggedright\arraybackslash}X}@{}}
    \toprule
    \textbf{ID} & \textbf{Riesgo} & \textbf{Probabilidad}
      & \textbf{Impacto} & \textbf{Mitigación} \\
    \midrule
    \pendiente{R-001} & \pendiente{descripción del riesgo} &
      \pendiente{Media} & \pendiente{Alto} &
      \pendiente{medida de mitigación} \\
    \addlinespace
    \pendiente{R-002} & \pendiente{descripción del riesgo} &
      \pendiente{Baja} & \pendiente{Medio} &
      \pendiente{medida de mitigación} \\
    \bottomrule
  \end{tabularx}
\end{table}
